Co-authorship is not a technique. It is what an interaction looks like 
when both people are genuinely present and neither is resolving meaning 
before it has formed between them.

It is recognisable not by what either person does, but by what the 
interaction produces. Meaning does not arrive from one side and land 
on the other. It accumulates in the space between --- built by both, 
owned by neither, carrying the fingerprints of each.

This is what distinguishes it from assertion with flipped dynamics. 
When only one person leaves meaning open and the other simply fills it, 
the structure has not changed --- only the direction. One person is still 
operating on the other. Co-authorship requires that both people are 
simultaneously incomplete, simultaneously reaching. Neither is the 
architect. Neither is the audience.

What this looks like varies. It can look like a silence neither person 
moves to fill, because neither needs to. It can look like a thought 
completed by the wrong person --- the one who didn't start it finishing 
it more accurately than the one who did. It can look like physical 
presence that neither initiated and neither interrupted. It can look 
like two people arriving at the same unspoken conclusion without 
either having directed the other there.

The common thread is not a behaviour. It is an absence: the absence 
of either person managing the interaction toward an outcome. When that 
absence is genuine on both sides, meaning builds on its own. 

Co-authorship is the symptom. Generative Incompletion, operating 
from both sides, is the condition that produces it.